\chapter{AI Ethics \& Bias Statement}
\label{app:ethics}

This appendix addresses the ethical considerations and potential biases associated with the multi-paradigm fraud detection system developed in this thesis, encompassing the Graph Neural Network (GNN), Spiking Neural Network (SNN), and Autoencoder (AE) ensemble pipelines. The deployment of AI-based fraud detection in financial services carries significant ethical implications, as automated decisions can directly affect individuals' access to credit and financial services. A thorough ethical analysis is therefore not merely an academic exercise but a professional obligation for any researcher or practitioner working in this domain~\cite{chouldechova2017fair}.

\section{Ethical Considerations}

The deployment of AI-based fraud detection systems raises several ethical concerns that must be carefully addressed across all three pipelines:

\textbf{Transparency and Explainability:} All three pipeline architectures present distinct explainability challenges. The GNN pipeline's decision-making process is mediated by multiple layers of message passing, making it difficult to trace which relational patterns contributed to a specific fraud prediction. The Split GNN's edge classifier provides some degree of interpretability by identifying which edges the model considers indicative of same-class or different-class relationships, but the subsequent spectral filtering layers remain opaque. The SNN pipeline is perhaps the least interpretable of the three: the rate coding transformation, temporal spike dynamics, and ensemble distillation process create a complex, non-linear mapping from input features to predictions that resists straightforward explanation. The AE pipeline offers moderate interpretability through reconstruction error analysis---analysts can examine which features contributed most to the high reconstruction error that flagged an application---but the latent space representations themselves are not directly interpretable. Full explainability across all pipelines remains an open challenge that requires dedicated research into pipeline-specific explanation methods.

\textbf{Privacy:} The GNN pipeline's graph construction process captures relational patterns between applicants (shared devices, overlapping behavioral profiles), which raises particular privacy concerns. While identifying fields (email, phone, device ID, IP address) are excluded from user features, the graph topology itself may encode sensitive information: an adversary with partial knowledge of the graph structure could potentially infer attributes of other users through link analysis. The SNN and AE pipelines operate on individual feature vectors rather than relational data, which reduces but does not eliminate privacy risks. Organizations must ensure compliance with data protection regulations (GDPR, CCPA) when constructing and storing both the graph data used by the GNN pipeline and the tabular data used by the SNN and AE pipelines. Differential privacy mechanisms and federated learning approaches could be explored to mitigate these risks in future work.

\textbf{Autonomy and Human Oversight:} Automated fraud detection systems should not make final decisions without human oversight, regardless of the underlying pipeline. The recommended deployment architecture includes a human-in-the-loop review process where flagged transactions are reviewed by trained fraud analysts before any action is taken. This is particularly important for the SNN pipeline, which achieves the highest TPR but also flags more applications for review, increasing the operational burden on fraud analysts. The AE pipeline's anomaly-based approach naturally produces a ranked list of suspicious applications with associated anomaly scores, which can be used to prioritize analyst review. All pipelines should be deployed with configurable decision thresholds that allow organizations to adjust the trade-off between fraud detection sensitivity and operational workload.

\textbf{Environmental Impact:} The computational cost of training all three pipelines is substantial (estimated 400+ GPU-hours for the full experimental suite). The environmental impact of this computation, measured in terms of energy consumption and carbon emissions, should be acknowledged and minimized where possible. Researchers should consider using more efficient hyperparameter search strategies (e.g., Bayesian optimization instead of grid search) and sharing pre-trained models to avoid redundant computation.

\section{Bias Assessment}

\textbf{Demographic Bias:} The subgroup fairness analysis (Tables~\ref{tab:ch8-fairness-comparison} and~\ref{tab:ch8-fairness-employment} in Chapter~\ref{ch:comparison}) reveals variation in TPR across age groups and employment statuses for all three pipelines. The most concerning finding is the consistent under-detection of fraud in the 60--69 age group across all pipelines (TPR of 41.8--49.1\% depending on pipeline), compared to the 18--29 age group (TPR of 54.1--61.2\%). This age-based disparity is not unique to any single pipeline but reflects a systemic pattern in the data and model architectures. As Chouldechova~\cite{chouldechova2017fair} has demonstrated, such fairness trade-offs are often mathematically unavoidable when base rates differ across subgroups, making it imperative that practitioners are transparent about the specific compromises entailed by their choice of model and threshold. The SNN pipeline exhibits the smallest TPR range (12.1\% by age), suggesting that its temporal spike dynamics and ensemble strategy provide somewhat more equitable detection, but the disparity remains significant.

\textbf{Data Representation Bias:} The training data may underrepresent certain demographic groups or geographic regions, leading to poorer model performance for underrepresented populations. The cross-version evaluation highlights how distributional shifts can disproportionately affect certain subgroups. The BAF dataset's multi-version design reveals that performance on underrepresented subgroups degrades more sharply under distributional shift than performance on well-represented subgroups, compounding the fairness concern~\cite{dai2021fair}. All three pipelines inherit this data representation bias, though the AE pipeline's unsupervised pre-training on the legitimate class manifold may partially mitigate it by learning representations that are not directly conditioned on potentially biased fraud labels.

\textbf{Historical Bias:} The fraud labels used for training are generated by existing detection systems, which may themselves contain biases. If the labeling system disproportionately flags certain demographic groups, the model will learn and amplify these biases. This concern applies to all three supervised pipelines (GNN and SNN) and to the stacking meta-learner in the AE pipeline. The AE pipeline's unsupervised pre-training phase is less susceptible to historical bias because it does not use fraud labels, but the stacking phase reintroduces label dependency. Adversarial debiasing during training and equalized odds post-processing are recommended mitigation strategies across all pipelines~\cite{kozodoi2022fairness}.

\textbf{Pipeline-Specific Bias Concerns:} The GNN pipeline introduces a unique bias risk through the graph construction process: the choice of which features become shared attribute nodes determines which relational patterns are detectable. If the selected features correlate with protected attributes (e.g., housing status correlating with socioeconomic status), the graph topology may encode discriminatory patterns that the GNN amplifies through message passing. The SNN pipeline's rate coding introduces a different bias vector: the mapping from continuous features to spike firing rates creates a monotonic transformation that may disproportionately emphasize features with high dynamic range, potentially amplifying biases in those features. The AE pipeline's target manifold isolation strategy, while partially mitigating historical bias, introduces the risk of defining the ``legitimate'' manifold in a way that systematically excludes certain demographic groups whose application patterns differ from the majority.

\textbf{Mitigation Strategies:} Recommended mitigation strategies across all three pipelines include: (1)~regular fairness audits using disaggregated evaluation metrics, with specific attention to the 60--69 age group that shows consistent under-detection; (2)~adversarial debiasing during training to learn representations that are invariant to protected attributes, applied to all supervised training phases~\cite{kozodoi2022fairness}; (3)~equalized odds post-processing to equalize TPR and FPR across demographic groups, with pipeline-specific calibration; (4)~diverse training data collection to ensure representation of all relevant subgroups, particularly for the under-detected 60--69 age group; (5)~graph schema auditing for the GNN pipeline to ensure that the selected attribute nodes do not disproportionately encode protected attribute information; (6)~rate coding calibration for the SNN pipeline to ensure that the spike firing rate mapping does not introduce feature-level biases that disproportionately affect certain demographic groups; and (7)~reinforcement-based fairness optimization~\cite{nasif2026reinforcement}, which offers a promising avenue for dynamically adjusting decision boundaries to satisfy fairness constraints during deployment, adapting to distributional shifts that may exacerbate bias over time.
